\section{Step 29 Record the bit length and append the marker}\label{sha-s29}

Read the original byte total and shift it left by three with \texttt{dsll}. Keep that full bit count in \texttt{s3} while writing the marker at the next free buffer position. Advance the working used count after the store. The context's original-byte total remains unchanged: the marker belongs to padding rather than the message.

\subsection{Write and inspect}
Replace only the temporary instructions between \texttt{\# BEGIN fin\_marker} and \texttt{\# END fin\_marker}. Keep the supplied labels outside that region.

\Needspace{10\baselineskip}
\MARSCodeLabel{MIPS assembly}
\begin{minted}[samepage=true]{mars-idea}
    ld $s3, 32($s0)
    dsll $s3, $s3, 3
    daddu $t0, $s0, $s2
    li $t1, 128
    sb $t1, 48($t0)
    addiu $s2, $s2, 1
\end{minted}

Use recipe \texttt{s29-1} with the cumulative program through this edit. The exact recipe names the entry, prepared input and checkpoint. Observe the checkpoint before continuing to its temporary or normal return.

\subsection{Expected observation}
For abc, \texttt{s3} is twenty-four, the marker occupies buffer index three and the working used count is four.

\Needspace{9\baselineskip}
\subsection{Check your understanding}
Does appending 80 increase the original message length?

\begin{itemize}
  \item \textbf{A} Yes, by one byte.
  \item \textbf{B} Yes, by eight bytes.
  \item \textbf{C} No; the recorded length excludes padding.
\end{itemize}
